\documentclass[../Main.tex]{subfiles}

\begin{document}

\section{Grundlagen und Methodik} \label{sec:Grundlagen}
Im folgenden Kapitel finden sich Definitionen und Informationen, die zum Verständnis unserer Arbeit beitragen. Zuerst grenzen wir das Hack-and-Slash-Genre ab (Abschnitt~\ref{sec:Genredefinition}) und führen das MDA-Framework sowie ein (Abschnitt~\ref{sec:MDA}). Danach beschreiben und begründen wir unser Vorgehen in den einzelnen Arbeitsschritten, in der Reihenfolge, in der wir gearbeitet haben: den MDA-Fragebogen für die Referenzspielanalyse, die Leveldesign-Recherche, die Umsetzung und die Playtests. Fachbegriffe, die wir benutzen und die zum Verständnis unserer Texte nötig sind, werden zudem im Glossar im Anhang erklärt.

\begin{itemize}[leftmargin=1.5cm]
    \item \textcolor{red}{\textbf{Glossar} existiert noch nicht: im Anhang anlegen (z.\,B. Combo, I-Frames, Hitstun, Knockback, Encounter, Welle, Greybox, Spawn, Telegraphing) und danach hier mit \textbackslash ref darauf verweisen.}
\end{itemize}

\subsection{Definition und Einordnung des Genres} \label{sec:Genredefinition}
Ein Videospielgenre zu definieren ist nicht immer einfach, denn Genres sind nicht statisch, sondern entwickeln sich laufend weiter und verzweigen sich. Zudem existiert selten eine scharfe Definition; vielmehr ergibt sich ein Genre aus den Gemeinsamkeiten vieler Spiele, wobei nicht jedes Spiel alle Merkmale erfüllen muss. Im Folgenden beschreiben wir deshalb unser Verständnis des Hack-and-Slash-Genres als Arbeitsdefinition für diese Arbeit. In der Videospielindustrie bezeichnet \textit{Hack-and-Slash} ein Untergenre der Action-Spiele. Der Fokus liegt dabei auf Echtzeitkampf mit Nahkampfwaffen wie Schwertern, Äxten oder anderen Klingen. Fernkampfwaffen sind zwar vertreten, spielen aber nur eine untergeordnete und unterstützende Rolle. Das Genre entstammt den 2D-Beat-'em-up-Spielen, entwickelte sich in der 3D-Ära weiter und wurde vor allem durch Titel wie \game{Devil May Cry} (2001)\autocite{capcom_devilmaycry_2001} oder \game{Bayonetta} (2009)\autocite{platinumgames_bayonetta_2009} geprägt.\autocite{wikipedia_hackandslash}

Moderne Hack-and-Slash-Spiele zeichnen sich durch hohe Spielgeschwindigkeiten und stilvolle, fliessende Bewegungen aus. Das Verketten präziser Angriffs-Combos, das gezielte Abbrechen von Bewegungen, I-Frames und strategisches Distanzmanagement zum Gegner prägen das Gameplay. Was ein überzeugendes Hack-and-Slash ausmacht, ist dabei nicht allein die Anzahl dieser Mechaniken, sondern ihr Zusammenspiel. Erst wenn präzises Kampfverhalten und befriedigendes \textcolor{red}{Spielgefühl} aufeinander abgestimmt sind, entsteht das genredefinierende Gefühl von Kontrolle und Ausdrucksstärke.\autocite[5--7]{lindgren2025striking}

Innerhalb des Genres hat sich für Spiele, die diese Ausdrucksstärke ins Zentrum stellen, die Bezeichnung \textit{Stylish-Action} etabliert (von englisch \textit{stylish}, stilvoll): Titel wie die bereits genannten \game{Devil May Cry} und \game{Bayonetta} fordern den Spieler auf, nicht nur effektiv, sondern möglichst elegant und abwechslungsreich zu kämpfen, und bewerten dies teils explizit mit Stil-Wertungen im Kampf.

\game{Project-Polearm} ist ein \textit{lineares} Hack-and-Slash und legt den Fokus besonders auf \textit{Fast-paced Combat} (englisch für temporeichen Kampf), also auf schnelle Kämpfe mit wenigen Pausen. Die Spieler sollen für aggressives Verhalten belohnt werden (Abschnitt~\ref{sec:V1}). Dabei soll sich die Spielfigur nicht schwerfällig, sondern agil durch die Level der Spielwelt bewegen.

\subsection{Das MDA-Framework und Game Feel} \label{sec:MDA}
Um das Zusammenspiel von Gegnern und Spielfigur in Referenzspielen systematisch zu erfassen und vergleichbar zu machen, braucht es einen geeigneten Analyserahmen. Dafür führen wir in diesem Abschnitt das MDA-Framework als methodisches Werkzeug ein. Es beschreibt strukturiert die Beziehung zwischen Spielregeln, emergenten Verhaltensmustern und dem emotionalen Spielerleben.

Das MDA-Framework (Mechanics, Dynamics, Aesthetics) wurde von Hunicke, LeBlanc und Zubek entwickelt und ist ein formaler Ansatz zur Analyse und zum Verständnis von Gamedesign. Es beschreibt Spiele nicht als statisches Medium wie Buch oder Film, sondern als interaktives Erlebnis, das aus dem Zusammenspiel dreier klar abgegrenzter Ebenen entsteht.\autocite[\textcolor{red}{S.~??}]{hunicke2004mda}

Die unterste Ebene, die Mechaniken, umfasst die konkreten Regeln und Systeme eines Spiels. Darunter versteht man alles, was durch Daten und Algorithmen getrieben ist, wie Schadenswerte, Treffererkennung, Bewegungsgeschwindigkeit oder die Gegner-KI. Daraus entstehen zur Laufzeit die Dynamiken, also das emergente Verhalten, das sich aus dem Zusammenspiel von Mechaniken und Spielerentscheidungen ergibt. Konkret geht es etwa um den Kampfrhythmus, taktische Ausweichmuster oder eskalierende Gegnerwellen. Die oberste Ebene, die Ästhetik (Aesthetics), beschreibt schliesslich die emotionale Reaktion des Spielers auf dieses System. Bei Hack-and-Slash-Spielen stehen Gefühle wie Kontrolle und Befriedigung bei gelungenen Combos im Vordergrund.

Entscheidend ist die Kausalität zwischen den Ebenen: Mechaniken erzeugen Dynamiken, die wiederum Ästhetik erzeugen. Die Spieler erleben das Spiel jedoch in umgekehrter Reihenfolge.\autocite[\textcolor{red}{S.~??}]{hunicke2004mda} Für die Analyse bestehender Spiele bedeutet das, dass man beim Erleben (Aesthetics) beginnt und sich systematisch zu den Ursachen, also den Mechaniken, vorarbeitet.

Unserer Einschätzung nach eignet sich dieses Framework besonders gut für die Analyse von Hack-and-Slash-Spielen, da das Genre ein besonders hohes Mass dieser Kausalität aufweist. Beispielsweise sind die genauen Parameter einer einzelnen Mechanik, etwa einer Attacke, bedeutungslos, wenn sie nicht im Kontext weiterer Systeme wie Gegnerlebenspunkten, Gegnerverhalten und Animationsgeschwindigkeit betrachtet werden.

Dieses Verständnis von Spielen als Zusammenspiel mehrerer Ebenen weist Parallelen zu Steve Swinks Konzept des \textit{Game Feel} auf, das er in seinem Buch \textit{Game Feel: A Game Designer’s Guide to Virtual Sensation} beschreibt. Swink definiert Game Feel als die subjektive Wahrnehmung der Steuerung und Rückmeldung eines Spiels und untersucht, weshalb sich bestimmte Spielmechaniken für den Spieler präzise, befriedigend oder wirkungsvoll anfühlen.\autocite[1--6]{swink2008gamefeel} Wenn in dieser Arbeit von \textit{\textcolor{red}{Spielgefühl}} die Rede ist, insbesondere in der Erkenntnisfrage, ist damit Game Feel im Sinne Swinks gemeint; die beiden Begriffe werden im Folgenden synonym verwendet.

\begin{itemize}[leftmargin=1.5cm]
    \item \textcolor{red}{\textbf{Begriffe Spielgefühl und Kampfgefühl:} Laut deiner Angabe meint \enquote{Spielgefühl} im Bericht je nach Kontext das Gesamtgefühl des Spiels oder (meistens) das Kampfgefühl. Das widerspricht dem Satz oben, dass Spielgefühl und Game Feel synonym verwendet werden. Alle Vorkommen von \enquote{Spielgefühl} im Bericht sind rot markiert.}
    \item \textcolor{red}{\textbf{Hier die drei Qualitäten der Erkenntnisfrage definieren} (Spielfluss, Schwierigkeit, Spielgefühl bzw.\ Kampfgefühl), am besten als eigener kurzer Absatz oder Unterabschnitt \enquote{Begriffe der Erkenntnisfrage}. Abschnitt~\ref{sec:Erkenntnisfrage} verweist hierher. Vorschlag dazu im Chat.}
\end{itemize}

\subsection{MDA-Fragebogen} \label{sec:Analysefragebogen}

Auf Basis des MDA-Frameworks haben wir einen strukturierten Fragebogen mit 15 Fragen entwickelt, der die drei Ebenen getrennt betrachtet. Dieser MDA-Fragebogen ist unser Analyseframework für die Referenzspielanalyse in Kapitel~\ref{sec:Referenzspielanalyse}. Die 15 Fragen sind den drei MDA-Ebenen zugeordnet und wurden in umgekehrter Reihenfolge gestellt, da die Analyse beim Erleben (Aesthetics) beginnt und sich systematisch zu den Ursachen vorarbeitet. Thematisch sind die Fragen so ausgerichtet, dass sie gezielt Gegnerverhalten, Gegnerplatzierung, Spielfluss, Schwierigkeit und Feedback beleuchten, da diese Aspekte im Zentrum unserer Erkenntnisfrage stehen.

Die Spieletest wurden je eines von Zhile Zhu(Ninja Gaiden) und eins von Rodin Arslan(Nier) gespielt. Das es von verschiedenen playtestern gespielt wurde spielt dabei keine rolle, da das MDA-Framework so ausgelegt ist, dass es dies erlaubt.

\subsubsection{Aesthetics (Ästhetik und emotionales Erleben)}

\begin{enumerate}[label=\textbf{A\arabic*}, leftmargin=1.5cm, itemsep=0.5em]
    \item Welche emotionale Reaktion löst der Kampf primär aus: Kontrolle, Befriedigung, Anspannung oder Frustration? Warum?
    \item Wie gut vermittelt das audiovisuelle Feedback die Wucht und Präzision der Treffer?
    \item Vermittelt die Gegnerplatzierung ein Gefühl von intentionalem Design? Wirkt die Anordnung der Gegner durchdacht oder zufällig, und wie beeinflusst sie das Empfinden des Spielers?
    \item Verändert sich das emotionale Erleben im Verlauf des Levelabschnitts, und ist eine Spannungskurve erkennbar?
    \item Unterscheiden sich die Combos darin, welche Emotionen sie auslösen? Bringen sie den Spieler in Bedrängnis, oder entsteht ein Gefühl der Erleichterung?
\end{enumerate}

\subsubsection{Dynamics (Spieldynamik und Interaktion)}

\begin{enumerate}[label=\textbf{D\arabic*}, leftmargin=1.5cm, itemsep=0.5em]
    \item Wie ist der Levelabschnitt von Beginn bis Ende aufgebaut? Welche Muster sind erkennbar, und wie hängen sie mit der Spannungskurve zusammen?
    \item Erzeugt die Kombination verschiedener Gegnertypen innerhalb eines Encounters taktischen Druck, oder agieren die Gegner unabhängig voneinander?
    \item Entsteht durch Gegnerplatzierung, Spawnpunkte und Gegneranzahl ein flüssiger Spielfluss? Gibt es einen spürbaren Unterschied zwischen Einzel- und Gruppenkonfrontationen?
    \item Wie erholt sich der Spieler nach einem Fehler? Ist die Feedback-Schleife zwischen Fehler und Korrektur nachvollziehbar?
    \item Gibt es verschiedene Combos bzw.\ Kampfabläufe, und folgen daraus verschiedene Arten, die Gegner zu bekämpfen?
\end{enumerate}

\subsubsection{Mechanics (Spielmechaniken)}

\begin{enumerate}[label=\textbf{M\arabic*}, leftmargin=1.5cm, itemsep=0.5em]
    \item Wie unterscheiden sich die Gegnertypen auf Mechanikebene: Bewegungsgeschwindigkeit, Angriffsreichweite, Schadenswerte?
    \item Wie ist das Angriffsverhalten der Gegner implementiert? Reagieren sie aktiv auf den Spieler, oder folgen sie festen, vordefinierten Mustern?
    \item Wie begrenzt das Leveldesign die Bewegungsoptionen von Spieler und Gegner? Gibt es räumliche Einschränkungen, die den Kampf beeinflussen?
    \item Wie sind das Hitstun- und das Knockback-System gestaltet? Lassen sich verschiedene Arten beobachten?
    \item Wie hängen die gedrückten Tasten mit den visuellen Kampfereignissen zusammen? Steckt hinter der Tastenbelegung eine erkennbare Absicht?
\end{enumerate}

\begin{itemize}[leftmargin=1.5cm]
    \item \textcolor{red}{\textbf{Begründung des MDA-Fragebogens (Vertragspunkt):} Tabelle Frage (A1--M5) $\rightarrow$ abgedecktes Kriterium (Gegnerverhalten, Gegnerplatzierung, Spielfluss, Schwierigkeit, Feedback) $\rightarrow$ Begründung mit Quelle (z.\,B. Swink für Feedback, Totten/Rogers für Platzierung und Raum, GDC-Talks für Pacing, Hunicke für die MDA-Ebenen).}
    \item \textcolor{red}{Kurz beschreiben, wie die Analyse durchgeführt wurde: Plattform/Version, gespielte Abschnitte, ungefähre Spielzeit, wer gespielt hat, ob aufgezeichnet wurde.}
    \item \textcolor{red}{\textbf{Schwierigkeit:} Laut Text beleuchten die Fragen auch die Schwierigkeit, aber keine fragt direkt danach (am nächsten: A1 Frustration, D4 Fehlerkorrektur). In der Begründungstabelle zeigen, welche Fragen sie abdecken, oder als Limitation nennen. Keine neuen Fragen einfügen, die Analyse wurde mit diesen gemacht.}
    \item \textcolor{red}{\enquote{diente zudem als inhaltliche Basis für den Playtest-Fragebogen}: Stimmt das? Das Post-Questionnaire richtet sich nach V1--V8 (Abschnitt~\ref{sec:PostQuestionnaire}). Ggf.\ streichen oder präzisieren.}
    \item \textcolor{red}{Die Fragen habe ich nur sprachlich geglättet (Satzzeichen, Grammatik bei D1); inhaltlich sind sie unverändert.}
\end{itemize}

\subsection{Vorgehen bei der Leveldesign-Recherche} \label{sec:VorgehenLeveldesign}
Während die Referenzspielanalyse (Kapitel~\ref{sec:Referenzspielanalyse}) am konkreten Spiel untersuchte, wie das Gegnerverhalten \textcolor{red}{Spielgefühl}, Spielfluss und Schwierigkeit beeinflusst, stützte sich die Recherche zu Leveldesign und Gegnerplatzierung auf Fachliteratur und nicht auf eine systematische Analyse. Ziel war die gezielte Herleitung von Gestaltungsprinzipien für einen zusammenhängenden Kampfabschnitt im gewählten Genre.

Bei den Quellen handelt es sich vor allem um Fachliteratur und Erfahrungswerte von Entwicklern. Sie umfassen zwei Fachbücher die wir bei der Quellensuche im internet gefunden haben, \textit{Level Up! The Guide to Great Video Game Design} von Scott Rogers\autocite{rogers2014levelup} und \textit{An Architectural Approach to Level Design} von Christopher W. Totten\autocite{totten2014architectural}, sowie vier Fachvorträge der Game Developers Conference (GDC) die wir schon im vorhinein vor dieser arbeit kannten und wieder rausgesucht haben: \enquote{Ten Principles for Good Level Design}\autocite{gdc_tenprinciples2010}, \enquote{Level Design in a Day: A Series of First Steps}\autocite{gdc_leveldesignaday2020}, \enquote{Moving Mountains: Speedy Level Creation as a Desperate Indie Dev}\autocite{gdc_movingmountains2020} sowie \enquote{Level Design Workshop: Balancing Action and RPG in Horizon Zero Dawn Quests}\autocite{gdc_horizonquests2018}. Damit haben wir bewusst Buchtheorie und praktische Erfahrung verbunden.

Im Fokus standen besonders jene Themenbereiche, die für den Prototyp relevant sind. Da \game{Project-Polearm} nicht auf Exploration oder eine offene Spielwelt zielt, sondern kampfzentriert ist, konzentrierte sich die Recherche auf Themen wie Encounter-Design, räumliche Typologien, Pacing und Spannungskurve sowie Scaffolding. Die Quellen sind häufig allgemein gehalten und nicht auf ein bestimmtes Genre bezogen, weshalb wir sie in Bezug auf das Hack-and-Slash-Genre interpretiert haben. Auf diese Weise liessen sich in Kapitel~\ref{sec:Leveldesign} konkrete, praktisch umsetzbare Designprinzipien ableiten.

Es handelt sich überwiegend um Fach- und Praktikerwissen erfahrener Entwickler, nicht um empirische oder wissenschaftliche Studien. Zudem behandelt keine der Quellen das Hack-and-Slash-Genre unmittelbar. Diese genrespezifische Lücke wird durch die eigene Referenzspielanalyse in Kapitel~\ref{sec:Referenzspielanalyse} teilweise ausgeglichen.

\begin{itemize}[leftmargin=1.5cm]
    \item \textcolor{red}{\textbf{Auswahlkriterium} für die Recherchepunkte nennen: Sie wirken auf Gegnerverhalten oder Gegnerplatzierung, sind in einem linearen Kampfabschnitt umsetzbar und im Playtest prüfbar. Zeitmangel gehört dagegen in die Limitationen (Kap.~\ref{sec:Schlussfolgerungen}).}
    \item \textcolor{red}{\textbf{Quellenauswahl begründen:} Warum gerade diese zwei Bücher und vier Talks? Wie habt ihr sie gefunden (Empfehlung, GDC Vault, Bibliothek)? 1--2 Sätze.}
\end{itemize}

\subsection{Vorgehen bei der Umsetzung} \label{sec:VorgehenUmsetzung}
Aus den Recherchen haben wir acht Designvorschläge ausgearbeitet, die sich an der Erkenntnisfrage orientieren: je vier zum Gegnerverhalten und zur Gegnerplatzierung (Abschnitt~\ref{sec:Designvorschlaege}). Bei allen Vorschlägen haben wir uns an den drei Qualitäten der Erkenntnisfrage orientiert.

Die praktische Umsetzung erfolgte in mehreren Stufen. Das Level haben wir zunächst in TrenchBroom als Greybox aufgebaut und danach in Blender optisch aufgewertet. Animationen und Character-Rigs stammen von Mixamo und wurden von uns für unsere Zwecke angepasst. Zusätzlich haben wir ein Animationspaket von 9CG verwendet, dessen Animationen wir ausnahmslos in Blender abgeändert haben. Musik und weitere Assets stammen von Online-Quellen wie Kenney, Poly Haven und weiteren Webseiten und wurden ebenfalls an unsere Bedürfnisse angepasst. Anschliessend haben wir alles in der Game Engine Godot zusammengeführt, in der auch der Grossteil des Codes entstand.

Die genauen Raumgrössen, Schadenswerte und weitere Parameter haben wir iterativ selbst erarbeitet, da wir unserer Meinung nach genügend eigene Spielerfahrung mitbringen. Dabei hatten wir unsere Recherche stets im Hinterkopf. Systeme wie der Gegnermanager (Abschnitt~\ref{sec:V2}) wurden direkt aus der Recherche heraus entworfen und gezielt für unsere Designvorschläge gestaltet. Genaueres zu unserer Code-Architektur findet sich in der technischen Dokumentation.\footnote{\textcolor{red}{[Titel der technischen Dokumentation einfügen]}} natürlich haben wir auch git und github zur zusammenarbeit benutzt und für den versionsveraluf. itterativ hies bei uns zuerst die systeme implementieren und danach schauen ob sie sich gut anfühlen wenn nicht dann noch einmal itereiren und wenn sie sich gut anfühlen dann schmücken. Im anhang ist eine liste mit externen assets, und auch zu unserem github, wo wir den sourcecode haben werden. im code sind dann auch noch refernezen.

\begin{itemize}[leftmargin=1.5cm]
    \item \textcolor{red}{Git (Versionsverwaltung) fehlt noch in der Werkzeugliste, falls ihr es verwendet habt.}
    \item \textcolor{red}{\enquote{iterativ}: Den Ablauf kurz konkret machen (Greybox $\rightarrow$ Umsetzung in Godot $\rightarrow$ selbst testen $\rightarrow$ anpassen $\rightarrow$ Playtest) und auf Abschnitt~\ref{sec:Prototyping} verweisen. Das stützt auch die Begründung mit der eigenen Spielerfahrung.}
    \item \textcolor{red}{Auf die Liste der externen Assets im Anhang verweisen (sobald sie steht) und die Arbeitsjournale als Beleg für den Prozess erwähnen.}
\end{itemize}

\subsection{Vorgehen bei den Playtests} \label{sec:VorgehenPlaytest}

\subsubsection{Methode und Ablauf} \label{sec:MethodePlaytest}
Ziel der Playtests war es, die in Abschnitt~\ref{sec:Designvorschlaege} hergeleiteten Designvorschläge mit der tatsächlichen Spielerfahrung von Testpersonen abzugleichen. Methodisch stützten wir uns dabei auf das Buch \textit{Games User Research} von Anders Drachen, Pejman Mirza-Babaei und Lennart E. Nacke.\autocite{drachen2018gamesuserresearch} Bei Games User Research (GUR) liegt der Fokus auf der Spielerfahrung, was genau zu unseren Zielen passt.

Durch die Kombination verschiedener Methoden lassen sich verschiedene Perspektiven und Ebenen einer Spielerfahrung erfassen.\autocite[\textcolor{red}{S.~??}]{drachen2018gamesuserresearch} So zeigt die Beobachtung, was eine Testperson tatsächlich tut, ein Fragebogen erfasst ihr subjektives Empfinden, und ein semistrukturiertes Interview kann gezielt auf Auffälligkeiten im Spielverhalten oder Empfinden eingehen.

Unser Ablauf war wie folgt: Zuerst wurden die Testpersonen mit einem Pre-Questionnaire charakterisiert. Darauf folgte der eigentliche Playtest, den wir beobachtet und aufgezeichnet haben, um die Reaktionen und das Spielverhalten der Testpersonen zu dokumentieren. Zum Abschluss füllten die Testpersonen ein Post-Questionnaire aus, und gegebenenfalls folgte ein semistrukturiertes Interview wenn wir weitere fragen hatten, die sich im laufe des playtest gestellt haben und wir die noch beantwortet haben wollten, um bestimmte Aussagen oder Verhaltensweisen genauer zu untersuchen.

Fragebögen (Questionnaires) eignen sich unter anderem zur Erfassung von Spielercharakteristika sowie von subjektiven Einstellungen und Erfahrungen.\autocite[\textcolor{red}{S.~??}]{drachen2018gamesuserresearch} Deshalb haben wir sie für genau diese beiden Zwecke im Pre- und im Post-Questionnaire eingesetzt.

Während des Playtests wurden zwei Aspekte gleichzeitig erfasst. Zum einen lässt das beobachtete Spielverhalten Rückschlüsse darauf zu, ob eine Testperson beispielsweise verwirrt ist, wiederholt Fehler macht, auffällige Entscheidungen trifft oder Orientierungsschwierigkeiten hat. Dabei haben wir zunächst nicht auf die Ursache geschlossen, sondern nur beobachtet und dokumentiert. Zum anderen wurden die Testpersonen aufgefordert, spontane Gedanken und emotionale Reaktionen laut zu äussern. Dieses laute Denken (\textit{Think-aloud}) gibt Einblick in die Wahrnehmung der Spielenden.\autocite[\textcolor{red}{S.~??}]{drachen2018gamesuserresearch} Im Kontext von Hack-and-Slash-Spielen ist eine kontinuierliche Verbalisierung jedoch schwierig, da sehr vieles gleichzeitig passiert und die kognitive Belastung in anspruchsvollen Situationen hoch sein kann. Deshalb haben wir die Methode angepasst: Die Testpersonen wurden nicht aufgefordert, ihre Gedanken durchgehend zu verbalisieren, sondern vor allem spontane und starke Reaktionen zu äussern. So konnten wir unmittelbare Reaktionen einfangen und Hinweise auf die Gedanken und Beweggründe hinter ihren Aktionen gewinnen.

Das semistrukturierte, retrospektive Interview diente dazu, während des Spielens beobachtete Auffälligkeiten und relevante Situationen nachträglich zu validieren und deren Ursachen aus Sicht der Testperson zu untersuchen.\autocite[\textcolor{red}{S.~??}]{drachen2018gamesuserresearch}

\begin{itemize}[leftmargin=1.5cm]
    \item \textcolor{red}{\textbf{Ablauf konsistent:} Hier steht \enquote{gegebenenfalls ein semistrukturiertes Interview}, in Abschnitt~\ref{sec:Interview} klingt es, als fände es immer statt und nutze die Videoaufnahme. So beschreiben, wie es tatsächlich gemacht wurde (Wurde die Aufnahme mit den Testpersonen angeschaut? Auch bei Discord?). Abweichungen gehören nach Abschnitt~\ref{sec:Durchfuehrung}, aber die Methode sollte nichts beschreiben, das nie stattfand.}
\end{itemize}

\subsubsection{Auswahl der Testpersonen} \label{sec:AuswahlTP}
Bei der Auswahl der Testpersonen haben wir uns nicht auf unseren engsten Bekanntenkreis beschränkt, sondern weit darüber hinaus gesucht. Besonders wichtig war es, Spielerinnen und Spieler zu finden, denen das Genre von \game{Project-Polearm} zusagt und die Erfahrung damit haben. Darüber hinaus haben wir uns überlegt, welche Perspektiven uns eine potenzielle Testperson liefern könnte, und versucht, möglichst unterschiedliche Typen zu finden. Wir haben beispielsweise eingeschätzt, wie ausführlich oder präzise eine Testperson antworten kann und welche Einblicke und Referenzen sie mitbringt. So haben wir unter anderem einen befreundeten Entwickler aus Deutschland angefragt, von dem wir besonders fachkundige Antworten erwartet haben. Zudem haben wir Frauen wie Männer einbezogen.

Die einzige Einschränkung, die uns aufgefallen ist: Alle Testpersonen sind in ihren späten Teenagerjahren oder Anfang zwanzig. Zudem spielt keine von ihnen ausschliesslich Hack-and-Slash-Spiele. Des weiteren kannten wir die meisten Testpersonen persönlich oder um 1 bis 2 ecken, was wohlwollende antworten zufolge haben könnte, deshalb haben wir auch ausdrücklich gesagt, dass sie das nicht machen sollen aber man weiss ja nie.

Wir haben fünf Playtests durchgeführt, weil dies im Rahmen unserer Maturitätsarbeit und angesichts des Umfangs unserer Befragung das Maximum war und laut unserer Betreuungsperson eine gute Anzahl ist, um genügend Perspektiven auf das Spiel zu erhalten.

Die testpersonen waren sich bewusst dass sie gefilmt wurden und haben dem zugestimmt und wir haben die aufnahmen nach auswertung mit ihnen wieder gelöscht.

\begin{itemize}[leftmargin=1.5cm]
    \item \textcolor{red}{\textbf{Mögliche Verzerrung} kurz nennen: Die meisten Testpersonen kannten euch persönlich; das kann zu wohlwollenderen Antworten führen. Ein Satz genügt, Details in den Limitationen.}
    \item \textcolor{red}{Einverständnis für die Aufnahmen und die Anonymisierung der Protokolle im Anhang erwähnen.}
    \item \textcolor{red}{\enquote{angesichts des Umfangs unserer Befragung}: Du hattest \enquote{Qualität unserer Fragen} geschrieben. Ich habe es als Umfang interpretiert, bitte prüfen.}
\end{itemize}

\subsubsection{Pre-Questionnaire} \label{sec:PreQuestionnaire}

Das Pre-Questionnaire diente dazu, die Testpersonen zu charakterisieren, damit ihre Aussagen im Kontext gesehen und gewichtet werden können.

\begin{enumerate}[label=\textbf{P\arabic*}, leftmargin=1.5cm, itemsep=0.5em]
    \item Wie viele Stunden spielst du pro Woche Videospiele?
    \item Was für Videospielgenres spielst du besonders gerne?
    \item Was macht für dich einen guten Kampfabschnitt aus?
    \item Was stellst du dir unter einem Hack-and-Slash vor?
    \item Wie schätzt du dein eigenes Können in actionlastigen Spielen generell ein? (Skala 1--5: \enquote{Anfänger} bis \enquote{sehr erfahren})
    \item Bevorzugst du eher vorsichtiges, methodisches Spiel oder aggressives, schnelles Spiel?
    \item Was erwartest du dir von diesem Playtest, bevor du ihn spielst?
    \item Worauf reagierst du in Kämpfen erfahrungsgemäss am empfindlichsten: z.\,B. unfaire Treffer, Unübersichtlichkeit, zu leichte Gegner, zu schwere Gegner?
\end{enumerate}

\begin{itemize}[leftmargin=1.5cm]
    \item \textcolor{red}{1 Satz, wofür die Antworten genutzt werden: P1, P5, P6 für das Spielerprofil (Tabelle in Abschnitt~\ref{sec:Durchfuehrung}, Spielstil in der Diskussion); P3, P7, P8 für Erwartungen, die mit Q14 verglichen werden.}
\end{itemize}

\subsubsection{Post-Questionnaire} \label{sec:PostQuestionnaire}

Das Post-Questionnaire erfasste die subjektive Einschätzung der Testpersonen direkt im Anschluss an den Playtest, pro Designvorschlag sowie für den gesamten Kampfabschnitt. Hinter jeder Aussage steht in Klammern der Designvorschlag, den sie überprüft (V1--V8, hergeleitet in Abschnitt~\ref{sec:Designvorschlaege}). Die Aussagen wurden auf einer Skala von 1 (\enquote{trifft gar nicht zu}) bis 5 (\enquote{trifft völlig zu}) bewertet, die offenen Fragen frei beantwortet.

\paragraph{Gegnerverhalten}
\begin{enumerate}[label=\textbf{Q\arabic*}, leftmargin=1.5cm, itemsep=0.5em]
    \item Ich fühlte mich meinen Gegnern grundsätzlich überlegen, ausser in einzelnen zugespitzten Momenten. \hfill (V1)
    \item Auch bei vielen Gegnern gleichzeitig behielt ich die Kontrolle über die Situation. \hfill (V2)
    \item Ich musste bewusst entscheiden, welchen Gegner ich zuerst angreife. \hfill (V3)
    \item Angriffe der Gegner wurden klar angekündigt, bevor sie mich trafen. \hfill (V4)
\end{enumerate}

\paragraph{Gegnerplatzierung und Kampfraum}
\begin{enumerate}[label=\textbf{Q\arabic*}, leftmargin=1.5cm, itemsep=0.5em, resume]
    \item Die Gegner waren zu Beginn klar überschaubar positioniert, das Verhalten wurde danach unvorhersehbarer. \hfill (V5)
    \item Der enge Abschnitt fühlte sich beim Spielen anders an als der offene Abschnitt. \hfill (V6)
    \item Der Kampf hatte spürbare Phasen von Anspannung und Erholung. \hfill (V7)
    \item Neue Gegnertypen wurden mir zuerst einzeln, dann erst in Kombination vorgestellt. \hfill (V8)
\end{enumerate}

\paragraph{Gesamteindruck}
\begin{enumerate}[label=\textbf{Q\arabic*}, leftmargin=1.5cm, itemsep=0.5em, resume]
    \item Insgesamt war dieser Kampfabschnitt für mich befriedigend zu spielen. \hfill (\textcolor{red}{Spielgefühl})
    \item Der gesamte Kampfabschnitt fühlte sich zusammenhängend an, nicht wie eine Aneinanderreihung einzelner Momente. \hfill (Spielfluss)
    \item Der Schwierigkeitsgrad war insgesamt angemessen. \hfill (Schwierigkeit)
    \item Was hat sich am besten angefühlt? \hfill (offen)
    \item Was hat sich unfair oder unangenehm angefühlt? \hfill (offen)
    \item Wie war dieser Kampfabschnitt im Vergleich zu dem, was du dir vor dem Test erwartet hast? \hfill (offen)
\end{enumerate}

\begin{itemize}[leftmargin=1.5cm]
    \item \textcolor{red}{Klarstellen, dass die Testpersonen die Klammern (V1--V8 usw.) nicht gesehen haben -- falls das so war. Sonst könnte es die Antworten beeinflusst haben.}
    \item \textcolor{red}{\textbf{Bekannte Schwächen} hier nicht ändern (so wurde gefragt), aber in der Diskussion bzw.\ den Limitationen aufgreifen: Q5 enthält zwei Aussagen auf einmal; Q6 spricht von \emph{dem} engen und \emph{dem} offenen Abschnitt, der Prototyp hat aber mehrere Narrow- und Prospect-Räume (Tabelle~\ref{tab:Raeume}).}
\end{itemize}

\subsubsection{Beobachtungsraster} \label{sec:Beobachtungsraster}

Während des Playtests wurde das Verhalten der Testperson anhand der folgenden Kriterien dokumentiert, jeweils mit Zeitstempel, um die Beobachtungen später dem passenden Segment und dem retrospektiven Interview zuordnen zu können. Dabei wurde zunächst nur beobachtet und protokolliert, nicht interpretiert.

\paragraph{Gegnerverhalten}
\begin{enumerate}[label=\textbf{B\arabic*}, leftmargin=1.5cm, itemsep=0.5em]
    \item Verhält sich die Testperson in normalen Phasen risikofreudig und aggressiv (z.\,B. wenig Ausweichen, viel Angreifen) und spürbar vorsichtiger in den zugespitzten Momenten? \hfill (V1)
    \item Verhält sich die Testperson in Momenten mit vielen Gegnern anders (hektischer, mehr Fehler) als in ruhigeren Momenten? \hfill (V2)
    \item Welchen Gegnertyp greift die Testperson zuerst an, wenn Nah- und Fernkämpfer gleichzeitig anwesend sind? \hfill (V3)
    \item Reagiert die Testperson (z.\,B. durch Ausweichen), bevor ein Angriff trifft, oder erst danach bzw.\ gar nicht? \hfill (V4)
    \item Wiederholt die Testperson denselben Fehler mehrfach an derselben Stelle?
\end{enumerate}

\paragraph{Gegnerplatzierung und Kampfraum}
\begin{enumerate}[label=\textbf{B\arabic*}, leftmargin=1.5cm, itemsep=0.5em, resume]
    \item Zeigt die Testperson zu Beginn eines Encounters ein zielgerichtetes Vorgehen (z.\,B. gezieltes Anvisieren einzelner Gegner), und wird dieses Vorgehen unruhiger, sobald sich die Gegner umzingelnd verhalten? \hfill (V5)
    \item Unterscheidet sich das Bewegungsmuster im engen Abschnitt sichtbar vom offenen Abschnitt (mehr Ausweichen, mehr Distanzhalten)? \hfill (V6)
    \item Zeigt die Testperson in Erholungsphasen ein erkennbar anderes Verhalten (z.\,B. Innehalten, freies Erkunden) als während einer Welle? \hfill (V7)
    \item Wirkt die Testperson beim ersten Auftreten eines neuen Gegnertyps zögerlich oder unsicher, und verschwindet dieses Zögern bei späteren Wiederholungen desselben Typs? \hfill (V8)
\end{enumerate}

\paragraph{Allgemein}
\begin{enumerate}[label=\textbf{B\arabic*}, leftmargin=1.5cm, itemsep=0.5em, resume]
    \item Anzeichen von Verwirrung oder Orientierungslosigkeit (z.\,B. Stehenbleiben, Umherblicken ohne erkennbares Ziel)
    \item Spontane verbale Reaktionen (Wortlaut und Zeitstempel notieren)
    \item Sichtbare Frustration oder sichtbare Freude bzw.\ Erleichterung, mit Zeitstempel
    \item Auffällige, nicht vorhergesehene Entscheidungen (z.\,B. Ignorieren eines Gegnertyps, unerwartete Fluchtversuche)
    \item Wirkt der Übergang zwischen zwei Segmenten (z.\,B. Welle zu Erholung, eng zu offen) beim Spielverhalten fliessend, oder gibt es einen sichtbaren Bruch (Zögern, Pause, sichtbare Neuorientierung)? \hfill (Spielfluss)
\end{enumerate}

\begin{itemize}[leftmargin=1.5cm]
    \item \textcolor{red}{\textbf{Zeitstempel:} In den Protokollen im Anhang fehlen sie teilweise. Entweder ergänzen oder hier \enquote{wo möglich} schreiben.}
    \item \textcolor{red}{Bei den Remote-Tests über Discord war nur der Bildschirm sichtbar. Kriterien wie \enquote{sichtbare Frustration} (B12) waren dort nur eingeschränkt beobachtbar. In Abschnitt~\ref{sec:Durchfuehrung} bzw.\ den Limitationen nennen.}
\end{itemize}

\subsubsection{Semistrukturiertes Interview} \label{sec:Interview}

Das Interview fand im Anschluss an das Post-Questionnaire statt und nutzte die Videoaufnahme des Playtests, um gezielt auf die während der Beobachtung notierten Auffälligkeiten einzugehen. Die folgenden Fragen dienten als Leitfaden und wurden je nach den tatsächlich beobachteten Auffälligkeiten der jeweiligen Testperson angepasst und ergänzt. Die tatsächlich gestellten Fragen pro Testperson finden sich im Anhang (Abschnitt~\ref{sec:PlaytestProtokolle}).

\begin{enumerate}[label=\textbf{I\arabic*}, leftmargin=1.5cm, itemsep=0.5em]
    \item An den markierten Stellen: Was ist dir hier durch den Kopf gegangen?
    \item Du hast hier [Reaktion/Aussage] gezeigt bzw.\ gesagt. Kannst du beschreiben, warum?
    \item An dieser Stelle bist du wiederholt gestorben bzw.\ hast du gezögert. Was hat dir das Weiterkommen erschwert?
    \item Vergleiche diesen Moment mit jenem dort: Haben sich die beiden unterschiedlich angefühlt, und wenn ja, wie?
    \item Als hier mehrere Gegner gleichzeitig da waren: Wie hast du entschieden, wen du zuerst angreifst?
    \item Empfandest du diesen Treffer bzw.\ diesen Tod als eigenen Fehler oder als unfair? Warum?
    \item Gab es einen Punkt, an dem sich der Kampf für dich in einen anderen \enquote{Modus} verschoben hat (z.\,B. von entspannt zu angespannt)? Wo genau?
    \item Gibt es etwas, das du während des Spielens gedacht oder gefühlt hast, das wir bisher nicht angesprochen haben?
\end{enumerate}

\subsubsection{Auswertung} \label{sec:Auswertung}
Da playtest eine menge arbeit sind und wir daher nicht viele von ihnen machen konnten im zusammenhang unserer Maturitätsarbeit, haben wir uns vor allem auf Qualitative Auswertung stat quantitavier konzentriert. Das hat zur folge, das wir keine abschliessenden aussagen zu der Erkenntnisfrage ausstellen können, da die menge der Testpersonen dafür zu klein ist. Ohne hin ist das thema rund um videospiele unserer meinung nach viel zu verwoben um einzelne faktoren als hauptasschlagspunkt bennenen zu können. Was wir mit der auswertung erreicht haben ist einfach die grobe richtung zu verstehen und die gründe einzugrenzen. des weiteren haben wir als resultat weitere bereiche und teilbereiche ausfindig gemacht, die weitere vielversprechende in unseren augen begründungen beiten können. Spezifisch haben wir zu den designvorschlägen jeweils ein paar fragen gehabt, die auf diese gezielt haben und konnten so versuchen unsere designvorschläge auszuwerten. dabei werden so aussagen wie x von 5 Testpersonen fallen. diese aussagen kann man auch in unserem anhang nachprüfuen. Wir konnten generell einfach tendenzen feststellen. Es sidn uns auch wiedersprüche in den antworten der testpersonen aufgefallen und auch wiedersprüche in unseren beobachtungen. diese werden in in der diskussion aufgegriffen. Zu den wiedersprüchen in unserer beobachtung, es liegt daran, dass sich das Spielverhalten der spieler vortwährend geändert hat und auch teils von raum zu raum unterschiedlich war so könenn mehrere sachen notiert worden sein die auf den ersten blick widersprüchlich wirkein es in anbetracht eines playtests aber nicht sind.
\begin{itemize}[leftmargin=1.5cm]
    \item \textcolor{red}{Die Auswertung ist qualitativ und beschreibend: Wegen der kleinen Stichprobe (fünf Testpersonen) werden keine Kennwerte wie Mittelwert oder Median berechnet (Absprache mit dem Betreuer; Playtests dieser Grösse dienen dem Erkennen von Mustern, nicht statistischen Aussagen).}
    \item \textcolor{red}{Fragebogen, Beobachtung und Interview werden pro Designvorschlag (V1--V8) gemeinsam betrachtet; Häufigkeiten werden als \enquote{x von 5 Testpersonen} angegeben, Likert-Antworten als Tendenz beschrieben.}
    \item \textcolor{red}{Widersprüche zwischen den Quellen werden in der Diskussion aufgegriffen. 2--4 Sätze reichen.}
\end{itemize}

\end{document}
